% =============================================================================
% sections/05_rio_analysis.tex  -  Section 5: RIO Analysis (second iteration)
% Teams: PM + test team / Status: draft (Preliminary data ported into
%   tables/rio_table.tex; PM review still needed - see that file's header)
%
% SCORING (0-20 pts):
%   Quality of updates to the Preliminary qualitative assessment (0-15):
%     risks, descriptions, mitigation actions and status; opportunities and
%     issues mentioned. LOW SCORE if the majority of risks are unchanged
%     (= risk assessment not performed). (0-10)
%   Trend chart (0-5): risk index changes across phases
%     (Preliminary Report -> Final Report), style: ECSS-M-ST-80C Fig. 5-7.
%
% NOTES:
%   - Risk IDs MUST match the Preliminary Report (R-01..R-10, O-01,
%     I-01..I-03). Removals/additions highlighted and justified.
%   - Jury feedback: L/S scale TABLES in the methodology; track L/S for
%     Issues and Opportunities as well; make generic risks
%     subsystem-specific (e.g. R-01).
%   - Font size 8 pt is allowed in RIO tables only.
% =============================================================================

\section{RIO Analysis}
\label{sec:rio}

\subsection{Methodology}
\label{subsec:rio_method}

The team applies a project-specific qualitative method during weekly subsystem
reviews (mechanical, electrical, software, science). For each identified risk,
issue, or opportunity, the owner defines the root cause, operational impact,
and actions, then assigns \textbf{Likelihood} (L) and \textbf{Severity} (S) on
a 1--5 scale (\autoref{tab:rio_scales}). The \textbf{Risk Index} is
$\mathrm{RI} = L \times S$. Likelihood and severity are tracked for
opportunities and issues as well as for risks. Response strategies:
\textit{Avoid}, \textit{Mitigate}, \textit{Accept}, \textit{Transfer},
\textit{Escalate} for risks and issues; \textit{Exploit}, \textit{Share},
\textit{Enhance}, \textit{Accept}, \textit{Escalate} for opportunities.

\begin{table}[H]
    \centering
    \caption{Likelihood and Severity scales}
    \label{tab:rio_scales}
    \begin{tabular}{C{0.5cm} L{5.6cm} | C{0.5cm} L{6.6cm}}
        \toprule
        \multicolumn{2}{c|}{\textbf{Likelihood (L)}} & \multicolumn{2}{c}{\textbf{Severity (S)}} \\
        \midrule
        1 & Rare ($p < 10\,\%$)                & 1 & Negligible (no replanning) \\
        2 & Unlikely ($10\,\% \leq p < 30\,\%$) & 2 & Minor (local workaround) \\
        3 & Possible ($30\,\% \leq p < 50\,\%$) & 3 & Moderate (subsystem rework or test delay) \\
        4 & Likely ($50\,\% \leq p < 70\,\%$)   & 4 & Major (milestone slip or partial task loss) \\
        5 & Almost certain ($p \geq 70\,\%$)    & 5 & Catastrophic (safety/compliance failure or mission task loss) \\
        \bottomrule
    \end{tabular}
\end{table}

\subsection{Risk Trend Chart}
\label{subsec:rio_trend}

\todo[inline]{Phase 2/4: trend chart (drawn outside \LaTeX, exported PDF)
with RI values per phase: Preliminary $\to$ Final. Preliminary baseline:
R-01:16, R-02:12, R-03:12, R-04:10, R-05:9, R-06:9, R-07:9, R-08:8, R-09:6,
R-10:6.}

\subsection{Qualitative Risks, Opportunities, and Issues Assessment}
\label{subsec:rio_register}

The second-iteration register of risks, opportunities, and issues, carried
over from the Preliminary Report with matching IDs, is provided as
\autoref{tab:rio} in the landscape appendix at the end of this document.
